\begin{table}[htbp]
\caption{Construct Status in ARC 0.9.0}
\label{tab:construct-status}
\centering
\small
\setlength{\tabcolsep}{3pt}
\renewcommand{\arraystretch}{0.95}
\begin{threeparttable}
\begin{tabular}{p{0.15\linewidth}p{0.17\linewidth}p{0.29\linewidth}p{0.27\linewidth}}
\toprule
Construct & Current status & Role & Retention or promotion rule \\
\midrule
\(\Anx,\Avo\) & Core person-level reference dimensions & Established attachment dimensions; H2 tests whether they moderate channel-specific regulation & Retain as attachment descriptors; remove ARC-specific dynamic moderation if H2 does not generalize \\
\(\Sec\) & Prespecified candidate core & Positive secure-base representation tested against recovery & Promote only after incremental value beyond \(\Anx,\Avo\) on a preregistered primary recovery outcome, out-of-sample improvement, and replication \\
\(\MentAff,\Flex\) & Optional extensions & Candidate pathway-specific moderators & Retain only with discriminant and incremental validity for their assigned pathways \\
\(Z^{\mathrm{obs}}\) & External validation criterion & Independently measured adult disorganization & Never define from ARC dynamics alone; use to test incremental validity of H5 indicators \\
\(C\) & Separate psychometric program & Candidate relational-control propensity & Not part of ARC unless independently validated and theoretically justified in a later model \\
\bottomrule
\end{tabular}
\begin{tablenotes}[flushleft]
\footnotesize
\item Note. Status refers to the ARC model architecture, not to whether the underlying construct exists in the broader attachment literature.
\end{tablenotes}
\end{threeparttable}
\end{table}
